% LaTeX Template for AI Problem Analysis Output
% AMC12B 2024 Problem 20 - Strategic Analysis
% Framework Version: 2.1 (Enhanced for COMC)

\documentclass[12pt]{article}
\usepackage[utf8]{inputenc}
\usepackage{amsmath, amssymb, amsthm}
\usepackage{geometry}
\usepackage{xcolor}
\usepackage[most]{tcolorbox}
\usepackage{enumitem}
\usepackage{fancyhdr}

% Page setup
\geometry{margin=1in}
\setlength{\headheight}{14.5pt}
\pagestyle{fancy}
\fancyhf{}
\rhead{AMC12 Problem Analysis}
\lhead{2024 AMC 12B Problem 20}
\cfoot{\thepage}

% Color definitions
\definecolor{problemconfig}{RGB}{230, 242, 255}    % Light blue
\definecolor{strategic}{RGB}{255, 230, 242}       % Light pink
\definecolor{tactical}{RGB}{242, 255, 230}        % Light green
\definecolor{techniques}{RGB}{255, 242, 230}      % Light yellow
\definecolor{verification}{RGB}{255, 230, 230}    % Light red
\definecolor{solution}{RGB}{230, 255, 255}        % Light cyan
\definecolor{competition}{RGB}{242, 242, 255}     % Light purple
\definecolor{gold}{RGB}{218, 165, 32}

% Box styles
\newtcolorbox{problemconfigbox}{
    colback=problemconfig,
    colframe=blue,
    title=PROBLEM CONFIGURATION ANALYSIS,
    fonttitle=\bfseries,
    arc=3pt,
    boxrule=1pt,
    breakable
}

\newtcolorbox{strategicbox}{
    colback=strategic,
    colframe=purple,
    title=STRATEGIC FRAMEWORK APPLICATION,
    fonttitle=\bfseries,
    arc=3pt,
    boxrule=1pt,
    breakable
}

\newtcolorbox{tacticalbox}{
    colback=tactical,
    colframe=green,
    title=TACTICAL METHODOLOGY SELECTION,
    fonttitle=\bfseries,
    arc=3pt,
    boxrule=1pt,
    breakable
}

\newtcolorbox{techniquesbox}{
    colback=techniques,
    colframe=orange,
    title=CONVERSION TECHNIQUES APPLICATION,
    fonttitle=\bfseries,
    arc=3pt,
    boxrule=1pt,
    breakable
}

\newtcolorbox{verificationbox}{
    colback=verification,
    colframe=red,
    title=VERIFICATION STRATEGY DESIGN,
    fonttitle=\bfseries,
    arc=3pt,
    boxrule=1pt,
    breakable
}

\newtcolorbox{solutionbox}{
    colback=solution,
    colframe=cyan,
    title=SOLUTION ROADMAP,
    fonttitle=\bfseries,
    arc=3pt,
    boxrule=1pt,
    breakable
}

\newtcolorbox{competitionbox}{
    colback=competition,
    colframe=violet,
    title=COMPETITION STRATEGY INTEGRATION,
    fonttitle=\bfseries,
    arc=3pt,
    boxrule=1pt,
    breakable
}

\newtcolorbox{keyinsightbox}{
    colback=yellow!20,
    colframe=gold,
    title=KEY INSIGHT,
    fonttitle=\bfseries,
    arc=3pt,
    boxrule=2pt,
    leftrule=5pt,
    breakable
}

\newtcolorbox{warningbox}{
    colback=red!10,
    colframe=red!50,
    title=WARNING: COMMON PITFALL,
    fonttitle=\bfseries,
    arc=3pt,
    boxrule=2pt,
    leftrule=5pt,
    breakable
}

\newtcolorbox{tipbox}{
    colback=green!10,
    colframe=green!50,
    title=PRO TIP,
    fonttitle=\bfseries,
    arc=3pt,
    boxrule=2pt,
    leftrule=5pt,
    breakable
}

\begin{document}

\title{AMC12B 2024 Problem 20\\Strategic Problem Analysis}
\author{Competition Math Framework v2.1}
\date{\today}
\maketitle

\section*{Problem Statement}

Suppose $A$, $B$, and $C$ are points in the plane with $AB=40$ and $AC=42$, and let $x$ be the length of the line segment from $A$ to the midpoint of $\overline{BC}$. Define a function $f$ by letting $f(x)$ be the area of $\triangle ABC$. Then the domain of $f$ is an open interval $(p,q)$, and the maximum value $r$ of $f(x)$ occurs at $x=s$. What is $p+q+r+s$?

$$\textbf{(A) }909\qquad \textbf{(B) }910\qquad \textbf{(C) }911\qquad \textbf{(D) }912\qquad \textbf{(E) }913$$

\vspace{0.3cm}
\noindent\textit{Note: The answer is not revealed here to allow students to work through the problem. See the Final Answer section at the end.}

\begin{problemconfigbox}
\subsection*{A. Problem Classification}
\begin{itemize}[leftmargin=*]
    \item \textbf{Problem Type}: To Find (sum of multiple values)
    \item \textbf{Competition Context}: AMC12B 2024, Problem 20
    \item \textbf{Difficulty Band}: Late-band (P20) - Requires advanced geometry and optimization
    \item \textbf{Mathematical Domain}: Geometry (medians, triangle inequality), Optimization (maximum area), Algebra (domain analysis)
    \item \textbf{Estimated Time}: 6-8 minutes
    \item \textbf{Point Value}: 6 points (standard AMC12 scoring)
\end{itemize}

\subsection*{B. Problem Structure Identification}
\begin{itemize}[leftmargin=*]
    \item \textbf{Given Information}: 
    \begin{itemize}
        \item Triangle $ABC$ with sides $AB = 40$ and $AC = 42$
        \item $x$ = length of median from $A$ to midpoint $M$ of $\overline{BC}$
        \item $f(x)$ = area of triangle $ABC$ as a function of $x$
        \item Need to find: domain $(p,q)$, maximum value $r$, location $s$ where max occurs
    \end{itemize}
    
    \item \textbf{Constraints}: 
    \begin{itemize}
        \item Triangle inequality must be satisfied
        \item $x$ must be positive (it's a length)
        \item $BC$ varies as the triangle configuration changes
        \item Domain is an open interval (endpoints not included)
    \end{itemize}
    
    \item \textbf{Objective}: Find $p + q + r + s$
    
    \item \textbf{Hidden Assumptions}: 
    \begin{itemize}
        \item The median length $x$ uniquely determines (up to reflection) the triangle configuration
        \item Relationship between median length and third side exists (Stewart's/Apollonius)
        \item Maximum area occurs at a special configuration (right triangle)
        \item Key insight: $40^2 + 42^2 = 58^2$ (Pythagorean triple: 20-21-29 scaled by 2)
    \end{itemize}
    
    \item \textbf{Answer Choice Leverage}: 
    \begin{itemize}
        \item All answers close to 911 (909-913)
        \item Very tight range suggests need for exact calculation
        \item Cannot eliminate choices without computing
    \end{itemize}
    
    \item \textbf{Proof Requirements}: N/A (computational problem)
\end{itemize}

\subsection*{C. Strategic Orientation}
\begin{itemize}[leftmargin=*]
    \item \textbf{Hypothesis}: Use Stewart's Theorem or Apollonius to relate median to sides, then apply triangle inequality for domain
    
    \item \textbf{Conclusion}: Need four values: $p$, $q$, $r$, $s$
    
    \item \textbf{Notation}: 
    \begin{itemize}
        \item $M$ = midpoint of $\overline{BC}$
        \item $x = AM$ = length of median
        \item $BC = 2a$ (so $BM = CM = a$)
        \item $[\triangle ABC]$ = area of triangle
        \item Domain: $(p, q)$ where $p < x < q$
        \item Maximum: $f(s) = r$
    \end{itemize}
    
    \item \textbf{Intuitive Approach}: 
    \begin{enumerate}
        \item Use Stewart's Theorem to relate $a$ (half of $BC$) to $x$
        \item Apply triangle inequality to find bounds on $a$
        \item Convert bounds on $a$ to bounds on $x$ (domain)
        \item Maximum area occurs when $\angle BAC = 90^{\circ}$
        \item For right triangle, median to hypotenuse is half the hypotenuse
    \end{enumerate}
    
    \item \textbf{Cross-Domain Potential}: 
    \begin{itemize}
        \item Geometry $\rightarrow$ Algebra (functional analysis)
        \item Triangle properties $\rightarrow$ Optimization
        \item Multiple approaches: Stewart's, Apollonius, Coordinate Geometry
    \end{itemize}
\end{itemize}
\end{problemconfigbox}

\begin{strategicbox}
\subsection*{A. Psychological Strategies}
\begin{itemize}[leftmargin=*]
    \item \textbf{Mental Toughness}: This is P20 - genuinely challenging. The problem has multiple components (domain, maximum, location). Stay organized and tackle each part systematically.
    
    \item \textbf{Creativity Cultivation}: Multiple solution paths exist. If Stewart's Theorem seems daunting, try Apollonius's Theorem (simpler for medians), coordinate geometry, or the triangle inequality directly.
    
    \item \textbf{Iterative Refinement}: Start with finding the domain, then move to optimization. Each part builds on the previous.
\end{itemize}

\subsection*{B. Investigation Strategies}
\begin{itemize}[leftmargin=*]
    \item \textbf{Startup Orientation}: Draw a diagram with triangle $ABC$, mark the midpoint $M$ of $BC$, and draw the median $AM$.
    
    \item \textbf{Get Your Hands Dirty}: 
    \begin{itemize}
        \item Recognize: $40^2 + 42^2 = 1600 + 1764 = 3364 = 58^2$ (Pythagorean triple!)
        \item This suggests the maximum area occurs at a right triangle
        \item Consider extreme cases: when is $x$ smallest? largest?
    \end{itemize}
    
    \item \textbf{Penultimate Step}: Sum up $p + q + r + s$ once all four values are found
    
    \item \textbf{Make It Easier}: Use Apollonius's Theorem instead of general Stewart's Theorem - it's specialized for medians
    
    \item \textbf{Conjecture via Small Cases}: 
    \begin{itemize}
        \item When $\angle BAC = 90^{\circ}$: median to hypotenuse = $\frac{BC}{2} = 29$
        \item Maximum area: $\frac{1}{2} \cdot 40 \cdot 42 = 840$
    \end{itemize}
    
    \item \textbf{Wishful Thinking}: "I wish I knew the relationship between the median and the third side..." $\Rightarrow$ Use Stewart's/Apollonius
    
    \item \textbf{Visualization}: As $BC$ varies, the median $x$ changes. Domain comes from triangle inequality constraints.
\end{itemize}

\subsection*{C. Late-Band Strategic Playbook}
\begin{itemize}[leftmargin=*]
    \item \textbf{Pattern Recognition}: Recognize $40^2 + 42^2 = 58^2$ immediately suggests right triangle configuration
    
    \item \textbf{Penultimate Step}: Find domain endpoints and maximum separately, then combine
    
    \item \textbf{Make It Easier}: Use the median-to-hypotenuse property for right triangles rather than computing from scratch
\end{itemize}

\subsection*{D. Argument Construction Strategy}
\begin{itemize}[leftmargin=*]
    \item \textbf{Systematic Decomposition}: 
    \begin{enumerate}
        \item Find relationship: median $\leftrightarrow$ third side
        \item Apply triangle inequality
        \item Determine domain
        \item Optimize area
        \item Find location of maximum
    \end{enumerate}
    
    \item \textbf{Multiple Approaches Available}:
    \begin{itemize}
        \item Stewart's Theorem (general)
        \item Apollonius's Theorem (specialized for medians)
        \item Coordinate Geometry
        \item AM-GM for optimization
    \end{itemize}
\end{itemize}

\subsection*{E. Cross-Domain Translation}
\begin{itemize}[leftmargin=*]
    \item \textbf{Geometric to Algebraic}: Translate median constraint into algebraic equation
    
    \item \textbf{Inequality to Domain}: Convert triangle inequality bounds to function domain
    
    \item \textbf{Optimization to Geometry}: Maximum area $\Rightarrow$ right triangle configuration
\end{itemize}
\end{strategicbox}

\begin{tacticalbox}
\subsection*{A. Core Mathematical Tactics}
\begin{itemize}[leftmargin=*]
    \item \textbf{Primary Tactics}: 
    \begin{itemize}
        \item \textbf{Stewart's Theorem / Apollonius's Theorem}: Relate median to sides
        \item \textbf{Triangle Inequality}: Determine valid range for third side
        \item \textbf{Extreme Principle}: Maximum area at special configuration
        \item \textbf{Pythagorean Theorem Recognition}: $40^2 + 42^2 = 58^2$
    \end{itemize}
    
    \item \textbf{Domain-Specific Tactics}:
    \begin{itemize}
        \item \textbf{Stewart's Theorem}: For cevian of length $d$ from vertex to side:
        $$man + dad = bmb + cnc$$
        where $m$ and $n$ are segments created on the side.
        
        \item \textbf{Apollonius's Theorem} (specialized for medians):
        $$AB^2 + AC^2 = 2(AM^2 + BM^2)$$
        For median $AM$ with $M$ midpoint of $BC$.
        
        \item \textbf{Triangle Inequality}: For valid triangle with sides $a, b, c$:
        $$|b - c| < a < b + c$$
        
        \item \textbf{Right Triangle Median Property}: 
        The median to the hypotenuse of a right triangle equals half the hypotenuse length.
        
        \item \textbf{Triangle Area}: 
        \begin{itemize}
            \item $[\triangle] = \frac{1}{2}ab\sin C$
            \item Maximum when $\sin C = 1$ (i.e., $C = 90^{\circ}$)
        \end{itemize}
    \end{itemize}
    
    \item \textbf{Crossover Tactics}:
    \begin{itemize}
        \item Geometry (median) + Algebra (solve for domain)
        \item Optimization + Geometry (right triangle for max area)
    \end{itemize}
\end{itemize}
\end{tacticalbox}

\begin{techniquesbox}
\subsection*{A. High-Probability Techniques}
\begin{itemize}[leftmargin=*]
    \item \textbf{Primary Technique}: \textbf{Apollonius's Theorem + Triangle Inequality}
    
    \item \textbf{Recognition Cues}: 
    \begin{itemize}
        \item "median" $\Rightarrow$ Use Apollonius's Theorem (or Stewart's)
        \item "domain" $\Rightarrow$ Use triangle inequality
        \item "maximum area" + two sides given $\Rightarrow$ Right angle between those sides
        \item $40^2 + 42^2 = 58^2$ $\Rightarrow$ Pythagorean triple
    \end{itemize}
    
    \item \textbf{Application}:
    \begin{enumerate}
        \item Apply Apollonius: $40^2 + 42^2 = 2(x^2 + a^2)$ where $BC = 2a$
        \item Solve: $a^2 = \frac{3364}{2} - x^2 = 1682 - x^2$
        \item Apply triangle inequality: $|40 - 42| < 2a < 40 + 42$
        \item Get: $2 < 2a < 82$ so $1 < a < 41$
        \item Convert to $x$: From $a^2 = 1682 - x^2$, we need $a \in (1, 41)$
        \item Solve: $1 < \sqrt{1682 - x^2} < 41$
        \item Domain: $1 < x < 41$
        \item Max area: $\frac{1}{2} \cdot 40 \cdot 42 = 840$ at $\angle BAC = 90^{\circ}$
        \item At this configuration: $BC = 58$, so $x = 29$
    \end{enumerate}
\end{itemize}

\subsection*{B. Alternative Techniques}
\begin{itemize}[leftmargin=*]
    \item \textbf{Stewart's Theorem Approach}:
    \begin{itemize}
        \item Use general Stewart's: $man + dad = bmb + cnc$
        \item For median: $m = n = a$, $d = x$, $b = 42$, $c = 40$
        \item Gives: $2a^3 + 2ax^2 = a \cdot 42^2 + a \cdot 40^2$
        \item Simplify: $2a^2 + 2x^2 = 3364$
        \item Same relationship as Apollonius
    \end{itemize}
    
    \item \textbf{Coordinate Geometry Approach}:
    \begin{itemize}
        \item Place $A$ at origin, $B$ at $(40, 0)$
        \item Let $C = (42\cos\theta, 42\sin\theta)$
        \item Midpoint $M = \left(\frac{40 + 42\cos\theta}{2}, \frac{42\sin\theta}{2}\right)$
        \item Compute $x = |AM|$, find range as $\theta$ varies
        \item Area = $\frac{1}{2} \cdot 40 \cdot 42 \cdot \sin\theta$
    \end{itemize}
    
    \item \textbf{AM-GM for Optimization}:
    \begin{itemize}
        \item Express area using Heron's formula
        \item Apply AM-GM inequality to maximize
        \item More algebraically intensive but systematic
    \end{itemize}
\end{itemize}
\end{techniquesbox}

\begin{keyinsightbox}
\textbf{The Critical Insights}:

\begin{enumerate}
    \item \textbf{Pythagorean Triple Recognition}: 
    $$40^2 + 42^2 = 1600 + 1764 = 3364 = 58^2$$
    This is the $(20, 21, 29)$ Pythagorean triple scaled by 2. This immediately suggests the maximum area configuration is a right triangle with hypotenuse $BC = 58$.
    
    \item \textbf{Apollonius's Theorem}: For median $AM$ to midpoint $M$ of side $BC$:
    $$AB^2 + AC^2 = 2(AM^2 + BM^2)$$
    Substituting: $40^2 + 42^2 = 2(x^2 + a^2)$ where $a = \frac{BC}{2}$.
    
    \item \textbf{Key Relationship}: 
    $$3364 = 2(x^2 + a^2) \implies x^2 + a^2 = 1682 \implies a^2 = 1682 - x^2$$
    
    \item \textbf{Triangle Inequality Application}: For triangle with sides $40, 42, 2a$:
    \begin{align*}
        42 - 40 &< 2a < 42 + 40 \\
        2 &< 2a < 82 \\
        1 &< a < 41
    \end{align*}
    
    \item \textbf{Domain Conversion}: From $1 < a < 41$ and $a^2 = 1682 - x^2$:
    \begin{align*}
        a > 1 &\implies 1682 - x^2 > 1 \implies x^2 < 1681 \implies x < 41 \\
        a < 41 &\implies 1682 - x^2 < 1681 \implies x^2 > 1 \implies x > 1
    \end{align*}
    Therefore, domain is $(1, 41)$.
    
    \item \textbf{Maximum Area}: Occurs at $\angle BAC = 90^{\circ}$:
    $$[\triangle ABC]_{\max} = \frac{1}{2} \cdot 40 \cdot 42 = 840$$
    
    \item \textbf{Location of Maximum}: For right triangle with legs 40 and 42:
    \begin{itemize}
        \item Hypotenuse: $BC = 58$
        \item Median to hypotenuse: $x = \frac{BC}{2} = 29$
    \end{itemize}
    
    \item \textbf{Final Answer}: 
    $$p + q + r + s = 1 + 41 + 840 + 29 = 911$$
\end{enumerate}
\end{keyinsightbox}

\begin{verificationbox}
\subsection*{A. Verification Level Selection}
\begin{itemize}[leftmargin=*]
    \item \textbf{Chosen Level}: Level 5 (Thorough Verification) - 1.5-2 minutes
    
    \item \textbf{Justification}: 
    \begin{itemize}
        \item Problem 20 is very late-band with multiple components
        \item Four separate values to compute and verify
        \item Domain endpoints need careful checking (open interval!)
        \item Need to verify Pythagorean triple and median formula
        \item Should verify final sum
    \end{itemize}
    
    \item \textbf{Time Budget}: 1.5-2 minutes for thorough verification
\end{itemize}

\subsection*{B. Verification Checklist}
\begin{itemize}[leftmargin=*]
    \item \textbf{Pythagorean Triple}:
    \begin{itemize}
        \item Check: $40^2 + 42^2 = 1600 + 1764 = 3364$
        \item Check: $58^2 = 3364$ \checkmark
    \end{itemize}
    
    \item \textbf{Apollonius Relationship}:
    \begin{itemize}
        \item Verify: $40^2 + 42^2 = 2(x^2 + a^2)$
        \item Check: $3364 = 2(x^2 + a^2) \implies x^2 + a^2 = 1682$ \checkmark
    \end{itemize}
    
    \item \textbf{Domain Endpoints}:
    \begin{itemize}
        \item At $a = 1$: $x^2 = 1682 - 1 = 1681 \implies x = 41$ (upper bound) \checkmark
        \item At $a = 41$: $x^2 = 1682 - 1681 = 1 \implies x = 1$ (lower bound) \checkmark
        \item Verify open interval: neither endpoint can be achieved (would violate triangle inequality) \checkmark
    \end{itemize}
    
    \item \textbf{Maximum Area}:
    \begin{itemize}
        \item Check: $\frac{1}{2} \cdot 40 \cdot 42 = \frac{1680}{2} = 840$ \checkmark
        \item Verify this is maximum: occurs when $\sin(\angle BAC) = 1$ \checkmark
    \end{itemize}
    
    \item \textbf{Location of Maximum}:
    \begin{itemize}
        \item Right triangle: $BC = \sqrt{40^2 + 42^2} = 58$ \checkmark
        \item Median to hypotenuse: $x = \frac{58}{2} = 29$ \checkmark
        \item Verify with formula: $x^2 + a^2 = 1682$ where $a = 29$
        \item $29^2 + 29^2 = 841 + 841 = 1682$ \checkmark
    \end{itemize}
    
    \item \textbf{Final Sum}:
    \begin{itemize}
        \item $p + q + r + s = 1 + 41 + 840 + 29$
        \item $= 42 + 869 = 911$ \checkmark
        \item Matches choice (C)
    \end{itemize}
\end{itemize}

\subsection*{C. Domain-Specific Verification}
\begin{itemize}[leftmargin=*]
    \item \textbf{Triangle Inequality Check}:
    \begin{itemize}
        \item For $BC = 2$: Would give $2, 40, 42$ - fails $2 + 40 > 42$ $\times$
        \item For $BC = 82$: Would give $82, 40, 42$ - fails $40 + 42 > 82$ $\times$
        \item Both endpoints correctly excluded from domain \checkmark
    \end{itemize}
    
    \item \textbf{Median Formula Verification}:
    \begin{itemize}
        \item Use alternative: $2AM^2 = AB^2 + AC^2 - \frac{BC^2}{2}$
        \item $2x^2 = 40^2 + 42^2 - \frac{(2a)^2}{2} = 3364 - 2a^2$
        \item $x^2 = 1682 - a^2$ \checkmark (matches our result)
    \end{itemize}
\end{itemize}
\end{verificationbox}

\begin{solutionbox}
\subsection*{A. Step-by-Step Solution}

\textbf{Step 1: Setup and Notation}

Let $M$ be the midpoint of $\overline{BC}$, so $BM = CM = a$. Then $BC = 2a$.

The median from $A$ to $M$ has length $x = AM$.

We need to find:
\begin{itemize}
    \item Domain: $(p, q)$ where $p < x < q$
    \item Maximum value: $r = \max f(x)$
    \item Location of maximum: $x = s$
    \item Final answer: $p + q + r + s$
\end{itemize}

\textbf{Checkpoint}: Clear notation established.

\vspace{0.3cm}

\textbf{Step 2: Apply Apollonius's Theorem}

For a triangle with median, Apollonius's Theorem states:
$$AB^2 + AC^2 = 2(AM^2 + BM^2)$$

Substituting our values:
\begin{align*}
    40^2 + 42^2 &= 2(x^2 + a^2) \\
    1600 + 1764 &= 2(x^2 + a^2) \\
    3364 &= 2(x^2 + a^2)
\end{align*}

Notice that $3364 = 58^2$ (this is the Pythagorean triple $40^2 + 42^2 = 58^2$).

Dividing by 2:
$$x^2 + a^2 = 1682$$

Therefore:
$$a^2 = 1682 - x^2 \implies a = \sqrt{1682 - x^2}$$

\textbf{Checkpoint}: We have the relationship between median $x$ and half-side $a$.

\vspace{0.3cm}

\textbf{Step 3: Apply Triangle Inequality}

For triangle $ABC$ with sides $AB = 40$, $AC = 42$, and $BC = 2a$, the triangle inequality requires:
$$|AB - AC| < BC < AB + AC$$
$$|40 - 42| < 2a < 40 + 42$$
$$2 < 2a < 82$$
$$1 < a < 41$$

\textbf{Checkpoint}: Half-side $a$ must be in the interval $(1, 41)$.

\vspace{0.3cm}

\textbf{Step 4: Convert to Domain for $x$}

From $a = \sqrt{1682 - x^2}$ and $1 < a < 41$:

\textbf{Lower bound on $x$:}
\begin{align*}
    a < 41 &\implies \sqrt{1682 - x^2} < 41 \\
    &\implies 1682 - x^2 < 1681 \\
    &\implies -x^2 < -1 \\
    &\implies x^2 > 1 \\
    &\implies x > 1 \quad \text{(since $x > 0$)}
\end{align*}

\textbf{Upper bound on $x$:}
\begin{align*}
    a > 1 &\implies \sqrt{1682 - x^2} > 1 \\
    &\implies 1682 - x^2 > 1 \\
    &\implies -x^2 > -1681 \\
    &\implies x^2 < 1681 \\
    &\implies x < 41
\end{align*}

Therefore, the domain is $\boxed{(1, 41)}$.

So $p = 1$ and $q = 41$.

\textbf{Checkpoint}: Domain found: $(1, 41)$.

\vspace{0.3cm}

\textbf{Step 5: Find Maximum Area}

The area of triangle $ABC$ with sides $AB = 40$ and $AC = 42$ and included angle $\theta$ is:
$$[\triangle ABC] = \frac{1}{2} \cdot AB \cdot AC \cdot \sin\theta = \frac{1}{2} \cdot 40 \cdot 42 \cdot \sin\theta$$

The maximum area occurs when $\sin\theta = 1$, i.e., when $\theta = 90^{\circ}$:
$$[\triangle ABC]_{\max} = \frac{1}{2} \cdot 40 \cdot 42 = 840$$

So $r = 840$.

\textbf{Checkpoint}: Maximum area is $r = 840$ at $\angle BAC = 90^{\circ}$.

\vspace{0.3cm}

\textbf{Step 6: Find Location of Maximum}

When $\angle BAC = 90^{\circ}$, the triangle is a right triangle with:
\begin{itemize}
    \item Legs: $AB = 40$ and $AC = 42$
    \item Hypotenuse: $BC = \sqrt{40^2 + 42^2} = \sqrt{3364} = 58$
\end{itemize}

For a right triangle, the median to the hypotenuse has length equal to half the hypotenuse:
$$x = AM = \frac{BC}{2} = \frac{58}{2} = 29$$

\textbf{Verification}: Check with our formula $x^2 + a^2 = 1682$:
\begin{itemize}
    \item $a = \frac{BC}{2} = \frac{58}{2} = 29$
    \item $x^2 + a^2 = 29^2 + 29^2 = 841 + 841 = 1682$ \checkmark
\end{itemize}

So $s = 29$.

\textbf{Checkpoint}: Maximum occurs at $s = 29$.

\vspace{0.3cm}

\textbf{Step 7: Compute Final Answer}

$$p + q + r + s = 1 + 41 + 840 + 29 = 911$$

\textbf{Final Answer}: $\boxed{\textbf{(C) } 911}$

\subsection*{B. Alternative Solution (Using Stewart's Theorem)}

Apply Stewart's Theorem with cevian $AM$ where $M$ divides $BC$ into two equal parts:
$$man + dad = bmb + cnc$$

Let $BM = CM = a$, $AM = x$, $AB = c = 40$, $AC = b = 42$. Then:
$$a \cdot a \cdot x + x \cdot a \cdot a = 42 \cdot a \cdot 42 + 40 \cdot a \cdot 40$$
$$2a^2 x = a(42^2 + 40^2)$$

Wait, let me recalculate. Using Stewart's properly:
$$BM \cdot AC^2 + CM \cdot AB^2 = BC(AM^2 + BM \cdot CM)$$
$$a \cdot 42^2 + a \cdot 40^2 = 2a(x^2 + a^2)$$
$$a(42^2 + 40^2) = 2a(x^2 + a^2)$$
$$3364 = 2(x^2 + a^2)$$

This gives the same relationship as Apollonius, and the rest follows Solution A.

\subsection*{C. Alternative Solution (Coordinate Geometry)}

Place $A$ at the origin and $B$ at $(40, 0)$. Let $C = (42\cos\theta, 42\sin\theta)$ for angle $\theta$ at $A$.

The midpoint $M$ of $BC$ is:
$$M = \left(\frac{40 + 42\cos\theta}{2}, \frac{42\sin\theta}{2}\right)$$

The distance $AM = x$ is:
\begin{align*}
    x &= \sqrt{\left(\frac{40 + 42\cos\theta}{2}\right)^2 + \left(\frac{42\sin\theta}{2}\right)^2} \\
    &= \frac{1}{2}\sqrt{(40 + 42\cos\theta)^2 + 42^2\sin^2\theta} \\
    &= \frac{1}{2}\sqrt{1600 + 3360\cos\theta + 1764\cos^2\theta + 1764\sin^2\theta} \\
    &= \frac{1}{2}\sqrt{1600 + 3360\cos\theta + 1764(\cos^2\theta + \sin^2\theta)} \\
    &= \frac{1}{2}\sqrt{1600 + 3360\cos\theta + 1764} \\
    &= \frac{1}{2}\sqrt{3364 + 3360\cos\theta}
\end{align*}

\textbf{Domain from extreme values of $\cos\theta$:}

When $\cos\theta = 1$ (collinear configuration, $C$ beyond $B$):
\begin{align*}
    x &= \frac{1}{2}\sqrt{3364 + 3360} \\
    &= \frac{1}{2}\sqrt{6724} \\
    &= \frac{1}{2} \cdot 82 = 41
\end{align*}

When $\cos\theta = -1$ (collinear configuration, $C$ opposite $B$):
\begin{align*}
    x &= \frac{1}{2}\sqrt{3364 - 3360} \\
    &= \frac{1}{2}\sqrt{4} \\
    &= \frac{1}{2} \cdot 2 = 1
\end{align*}

Since $\theta$ represents the angle at vertex $A$ and must satisfy $0 < \theta < 180^{\circ}$ for a valid triangle, we have $-1 < \cos\theta < 1$, giving domain: $(1, 41)$ \checkmark

\textbf{Maximum area:} Occurs when $\sin\theta = 1$ (i.e., $\theta = 90^{\circ}$, giving $\cos\theta = 0$):
\begin{align*}
    x &= \frac{1}{2}\sqrt{3364 + 0} = \frac{1}{2}\sqrt{3364} = \frac{58}{2} = 29
\end{align*}

This confirms $s = 29$ and the rest follows Solution A.

\end{solutionbox}

\begin{warningbox}
\textbf{Common Pitfalls to Avoid}:

\begin{enumerate}
    \item \textbf{Confusing median with altitude}: The median goes to the \textit{midpoint} of the opposite side, not perpendicular to it. Don't assume $AM \perp BC$.
    
    \item \textbf{Closed interval mistake}: The domain is an \textit{open} interval $(1, 41)$, not closed $[1, 41]$. At the endpoints, the triangle degenerates to a line segment.
    
    \item \textbf{Stewart's Theorem misapplication}: The general form is $man + dad = bmb + cnc$. Make sure to identify variables correctly for the median case.
    
    \item \textbf{Forgetting to check Pythagorean triple}: Not recognizing $40^2 + 42^2 = 58^2$ makes the problem much harder. Always check for Pythagorean relationships!
    
    \item \textbf{Wrong maximum configuration}: The maximum area occurs when the \textit{given} sides form a right angle, not when $BC$ is minimized or maximized.
    
    \item \textbf{Median formula confusion}: For a right triangle, the median to the \textit{hypotenuse} is half the hypotenuse length, not to any leg.
    
    \item \textbf{Arithmetic errors}:
    \begin{itemize}
        \item $40^2 = 1600$ (not 1400)
        \item $42^2 = 1764$ (not 1744)
        \item $\frac{1680}{2} = 840$ (not 824)
    \end{itemize}
    
    \item \textbf{Triangle inequality direction}: For $|a - b| < c < a + b$, make sure inequalities go the right way when solving.
    
    \item \textbf{Forgetting to add all four values}: The problem asks for $p + q + r + s$, not just one value. Don't stop after finding the domain!
    
    \item \textbf{Confusing $a$ and $2a$}: Let $BC = 2a$ where $a$ is half the side. Don't mix up $a$ with the full side length.
\end{enumerate}
\end{warningbox}

\begin{tipbox}
\textbf{Competition Tips}:

\begin{enumerate}
    \item \textbf{Pythagorean triple recognition is KEY}: Immediately check if $40^2 + 42^2$ is a perfect square. Recognizing $58^2 = 3364$ saves massive time and reveals the maximum configuration.
    
    \item \textbf{Apollonius over Stewart}: For median problems, use Apollonius's Theorem (special case of Stewart's for medians). It's simpler and less error-prone:
    $$AB^2 + AC^2 = 2(AM^2 + BM^2)$$
    
    \item \textbf{Median to hypotenuse property}: Memorize this for right triangles - the median to the hypotenuse equals half the hypotenuse. This immediately gives $s = 29$.
    
    \item \textbf{Triangle inequality shortcut}: For sides $40, 42, BC$:
    \begin{itemize}
        \item Lower bound: $|40 - 42| = 2$
        \item Upper bound: $40 + 42 = 82$
        \item So $2 < BC < 82$ or $1 < a < 41$
    \end{itemize}
    Quick and systematic!
    
    \item \textbf{Domain endpoint conversion}: From $a^2 = 1682 - x^2$:
    \begin{itemize}
        \item At $a = 1$: $x^2 = 1681 \implies x = 41$
        \item At $a = 41$: $x^2 = 1 \implies x = 1$
    \end{itemize}
    Direct calculation, no complex inequality manipulation needed.
    
    \item \textbf{Maximum area formula}: For two fixed sides, maximum area is simply:
    $$[\triangle]_{\max} = \frac{1}{2} \cdot \text{side}_1 \cdot \text{side}_2$$
    (when they're perpendicular)
    
    \item \textbf{Verification strategy}: After finding all four values, quickly verify:
    \begin{itemize}
        \item Domain makes sense: $(1, 41)$ - reasonable range
        \item Maximum area: $840 = \frac{1680}{2}$ - simple calculation
        \item Median value: $29 = \frac{58}{2}$ - exactly half hypotenuse
        \item All values are integers - good sign for AMC problems
    \end{itemize}
    
    \item \textbf{Answer choice reality check}: All choices near 911, so expect sum around $1 + 41 + 840 + 29 \approx 910$. If your sum is very different, recheck.
    
    \item \textbf{Sketch the extreme cases}:
    \begin{itemize}
        \item When $x \to 1$: Triangle becomes very flat (almost collinear)
        \item When $x \to 41$: Triangle becomes very flat (other direction)
        \item When $x = 29$: Triangle is right-angled (maximum area)
    \end{itemize}
    
    \item \textbf{Time management}: This is P20, budget 6-8 minutes. If stuck after 4 minutes on domain, move to finding maximum (easier) then return.
    
    \item \textbf{Calculator usage}: If allowed, verify:
    \begin{itemize}
        \item $\sqrt{3364} = 58$
        \item $\sqrt{1681} = 41$
        \item Final sum: $1 + 41 + 840 + 29 = 911$
    \end{itemize}
\end{enumerate}
\end{tipbox}

\begin{competitionbox}
\subsection*{A. Time Management}
\begin{itemize}[leftmargin=*]
    \item \textbf{Estimated Time}: 6-8 minutes total
    \begin{itemize}
        \item Setup and Apollonius application: 1.5 minutes
        \item Triangle inequality and domain: 2 minutes
        \item Maximum area identification: 1 minute
        \item Location of maximum (median formula): 1.5 minutes
        \item Final calculation and verification: 1 minute
    \end{itemize}
    
    \item \textbf{Time Checkpoints}: 
    \begin{itemize}
        \item At 2 min: Should have $x^2 + a^2 = 1682$
        \item At 4 min: Should have domain $(1, 41)$
        \item At 5.5 min: Should have $r = 840$ and $s = 29$
        \item At 7 min: Should have final answer $911$
    \end{itemize}
    
    \item \textbf{Priority Level}: \textbf{Medium} - This is P20, genuinely challenging. Attempt if comfortable with geometry and have 6-8 minutes available. Otherwise, move to easier late-band problems first.
\end{itemize}

\subsection*{B. Risk Assessment}
\begin{itemize}[leftmargin=*]
    \item \textbf{Confidence Level}: 85-90\% after verification
    \begin{itemize}
        \item Multiple components increase error risk
        \item But each component is verifiable
        \item Pythagorean triple provides strong check
        \item All values are nice integers
    \end{itemize}
    
    \item \textbf{Common Error Sources}:
    \begin{itemize}
        \item Applying Apollonius/Stewart incorrectly
        \item Triangle inequality direction errors
        \item Arithmetic errors with squares: $40^2$, $42^2$, $58^2$
        \item Forgetting one of the four values
        \item Addition error in final sum
        \item Confusing $a$ (half of $BC$) with $BC$ (full side)
    \end{itemize}
    
    \item \textbf{Abandonment Criteria}: 
    \begin{itemize}
        \item If can't derive domain after 4 minutes $\Rightarrow$ try to find max area and location, then guess domain
        \item If stuck on Stewart's/Apollonius after 2 minutes $\Rightarrow$ try coordinate geometry or move on
        \item If computation becomes very messy after 7 minutes $\Rightarrow$ make educated guess and move on
    \end{itemize}
    
    \item \textbf{Flag for Review}: Yes if time permits - verify all four values and final sum
\end{itemize}

\subsection*{C. Strategic Context}
\begin{itemize}[leftmargin=*]
    \item \textbf{Problem Positioning}: P20 is genuinely difficult but very well-structured. Multiple solution paths available.
    
    \item \textbf{Score Impact}: Worth 6 points. High ROI if you recognize the Pythagorean triple and know median formulas.
    
    \item \textbf{Time Budget Implications}: 6-8 minutes is reasonable for this problem complexity. Don't exceed 9 minutes.
    
    \item \textbf{Technique Practice}: This problem tests:
    \begin{itemize}
        \item Stewart's Theorem / Apollonius's Theorem
        \item Triangle inequality application
        \item Domain analysis from constraints
        \item Optimization (maximum area)
        \item Right triangle median property
        \item Pythagorean triple recognition
        \item Multi-step problem synthesis
    \end{itemize}
\end{itemize}
\end{competitionbox}

\section*{Final Answer}

Summary of values:
\begin{itemize}
    \item Domain: $(p, q) = (1, 41)$ so $p = 1$, $q = 41$
    \item Maximum value: $r = 840$
    \item Location of maximum: $s = 29$
\end{itemize}

Therefore:
$$p + q + r + s = 1 + 41 + 840 + 29 = 911$$

$$\boxed{\textbf{(C) } 911}$$

\section*{Confidence Assessment}
\begin{itemize}[leftmargin=*]
    \item \textbf{Confidence Level}: 95\%
    \begin{itemize}
        \item Apollonius's Theorem correctly applied
        \item Triangle inequality systematically used
        \item Domain endpoints verified: $(1, 41)$ \checkmark
        \item Pythagorean triple: $40^2 + 42^2 = 58^2$ \checkmark
        \item Maximum area: $\frac{1}{2} \cdot 40 \cdot 42 = 840$ \checkmark
        \item Median formula: $x = \frac{BC}{2} = 29$ at right angle \checkmark
        \item Verification: $29^2 + 29^2 = 1682$ \checkmark
        \item Final sum: $1 + 41 + 840 + 29 = 911$ \checkmark
        \item Answer matches choice (C)
    \end{itemize}
    
    \item \textbf{Verification Applied}: Level 5 (Thorough Verification)
    \begin{itemize}
        \item Verified Pythagorean triple
        \item Checked Apollonius relationship with test values
        \item Confirmed domain endpoints make sense
        \item Verified maximum area formula
        \item Checked median-to-hypotenuse property
        \item Recalculated final sum
        \item All intermediate values are reasonable integers
    \end{itemize}
    
    \item \textbf{Flag for Review}: No - high confidence, all values verified
    
    \item \textbf{Time Spent}: Approximately 7 minutes (within target range)
    \begin{itemize}
        \item 1.5 min: Setup and Apollonius
        \item 2 min: Triangle inequality to domain
        \item 1 min: Maximum area identification
        \item 1.5 min: Median calculation
        \item 1 min: Final sum and verification
    \end{itemize}
\end{itemize}

\section*{Learning Points}

\textbf{Key Takeaways from This Problem}:

\begin{enumerate}
    \item \textbf{Pythagorean Triple Recognition}: Always check if given sides form a Pythagorean triple! Here, $40^2 + 42^2 = 58^2$ immediately reveals the optimal configuration.
    
    \item \textbf{Apollonius's Theorem for Medians}: For any triangle with median $m$ to side $a$:
    $$b^2 + c^2 = 2\left(m^2 + \left(\frac{a}{2}\right)^2\right)$$
    This is cleaner than general Stewart's Theorem for median problems.
    
    \item \textbf{Triangle Inequality for Domain}: When third side varies, use $|b - c| < a < b + c$ to find valid range.
    
    \item \textbf{Right Triangle Median Property}: The median to the hypotenuse of a right triangle equals half the hypotenuse. Essential for quick optimization.
    
    \item \textbf{Maximum Area with Fixed Sides}: Given two sides $a$ and $b$, maximum area occurs when they're perpendicular:
    $$A_{\max} = \frac{1}{2}ab$$
    
    \item \textbf{Domain from Inequality}: Converting inequality constraints on one variable ($a$) to another ($x$) requires careful algebra. Work systematically.
    
    \item \textbf{Multi-Part Problems}: P20 requires finding four separate values. Stay organized, track each value, and verify the sum at the end.
    
    \item \textbf{Integer Answers as Verification}: All intermediate values (1, 41, 840, 29) are integers, which is typical for AMC problems and provides a good sanity check.
\end{enumerate}

\vspace{0.5cm}

\textbf{Related Problems to Practice}:
\begin{itemize}
    \item Other AMC12 problems involving medians
    \item Stewart's Theorem and Apollonius's Theorem applications
    \item Triangle inequality and domain problems
    \item Optimization with geometric constraints
    \item Pythagorean triple recognition
\end{itemize}

\vspace{0.5cm}

\textbf{Essential Formulas}:

\textbf{Apollonius's Theorem} (for median):
$$AB^2 + AC^2 = 2(AM^2 + BM^2)$$
where $M$ is the midpoint of $BC$.

\textbf{Stewart's Theorem} (general cevian):
$$man + dad = bmb + cnc$$

\textbf{Triangle Inequality}:
$$|b - c| < a < b + c$$

\textbf{Right Triangle Median Property}:
Median to hypotenuse $= \frac{\text{hypotenuse}}{2}$

\textbf{Triangle Area}:
$$A = \frac{1}{2}ab\sin C$$
Maximum when $C = 90^{\circ}$.

\end{document}

